\documentclass[11pt]{article}

\usepackage[margin=1in]{geometry}
\usepackage{amsmath}
\usepackage{booktabs}
\usepackage[hidelinks]{hyperref}

\title{Power Profile of a Cyclist:\\ Pacing and Performance over Time-Trial Courses}
\author{Carter Dandridge \and Ben Miller \and Dennis Lee \and Prannay Veerabahu}
\date{\today}

\begin{document}
\maketitle

% ====================================================================
% NOTE: Archived LaTeX scaffold, updated to describe the CANONICAL
% time-stepping model (notebooks/model_v3.ipynb). The team's actual
% write-up is a shared document following docs/paper-outline.md; this
% file is kept only as a structural reference and is not maintained.
% ====================================================================

\section*{Summary}
% TODO (write yourselves, last): a standalone one-page overview.

\section{Introduction}

\subsection{Background}
% TODO (write yourselves): individual time trials, and why pacing, rider type, and terrain matter.

\subsection{Problem restatement}
% TODO (write yourselves): given a rider's power limits, how should power be distributed over a fixed course to minimise time, under accumulating fatigue?

\subsection{Variables and notation}
\begin{center}
\begin{tabular}{cl}
\toprule
Symbol & Meaning \\
\midrule
$v$ & speed (m/s) \\
$w$ & effective headwind on a segment (m/s) \\
$d_i$ & length of segment $i$ (m) \\
$\theta$ & slope angle, $\arctan(\mathrm{grade}\%/100)$ (rad) \\
$\tau_i$ & turn penalty on segment $i$ (s) \\
$\Delta t$ & time step ($1$ s) \\
$\rho$ & air density (kg/m$^3$) \\
$C_dA$ & drag area (m$^2$) \\
$C_{rr}$ & rolling-resistance coefficient \\
$m$ & rider $+$ bike mass (kg) \\
$g$ & gravitational acceleration (m/s$^2$) \\
$P_\mathrm{out}$ & power actually applied (W) \\
$P_\mathrm{base}$ & flat-ground target power (W) \\
$P_\mathrm{thr}$ & power threshold above which fatigue builds (W) \\
$P_\mathrm{max}$ & short-duration maximal power ceiling (W) \\
$k_\mathrm{hill}$ & gradient-to-power sensitivity (W/rad) \\
$f$ & fatigue state, $0 \le f \le 1$ \\
$r_\mathrm{fat},\, r_\mathrm{rec}$ & fatigue build / recovery rates \\
$c_\mathrm{fat}$ & fatigue reduction of max power \\
$v_\mathrm{max}$ & descent speed cap (m/s) \\
\bottomrule
\end{tabular}
\end{center}

\subsection{Assumptions}
% TODO (write yourselves): JUSTIFY each one -- the rubric grades the justification, not the list.
\begin{enumerate}
  \item The rider and bike are a single point mass.
  \item The course is a sequence of segments, each with constant gradient and wind exposure and a fixed turn penalty.
  \item Motion is integrated in discrete one-second steps (forward Euler); forces are held constant within a step.
  \item Target power rises linearly with gradient, $P_\mathrm{base} + k_\mathrm{hill}\theta$ (floored at $50$ W): the rider pushes harder uphill and eases downhill.
  \item Fatigue accumulates when output exceeds $P_\mathrm{thr}$ (growing with the squared relative overshoot) and recovers below it, stays in $[0,1]$, and scales down the available maximum power.
  \item Wind enters as an effective scalar headwind per segment, set by the segment's exposure label.
  \item Turns cost a fixed time penalty, not an explicit deceleration.
  \item Descent speed is capped at $v_\mathrm{max}$.
\end{enumerate}

\subsection{Modelling approach}
% TODO (write yourselves): why a time-stepping force balance with a gradient-driven power target and a threshold-fatigue mechanic is reasonable for this problem.

\section{The model}

The rider and bike are treated as one mass advanced in one-second steps. On a segment of slope angle $\theta = \arctan(\mathrm{grade}\%/100)$ the resisting forces are aerodynamic drag, rolling resistance, and gravity:
\begin{equation}
F_\mathrm{drag} = \tfrac12 \rho\, C_dA\, (v + w)^2, \qquad
F_\mathrm{roll} = C_{rr}\, m\, g \cos\theta, \qquad
F_\mathrm{grav} = m\, g \sin\theta .
\end{equation}
The rider applies power $P_\mathrm{out}$, giving propulsive force $F_\mathrm{rider} = P_\mathrm{out} / \max(v, 0.5)$ (the floor prevents division by zero at the start). Target power rises with gradient and is capped by a fatigue-reduced ceiling:
\begin{equation}
P_\mathrm{target} = \max\!\big(50,\ P_\mathrm{base} + k_\mathrm{hill}\theta\big), \qquad
P_\mathrm{out} = \min\!\big(P_\mathrm{target},\ P_\mathrm{max}(1 - c_\mathrm{fat} f)\big).
\end{equation}
Fatigue builds above the threshold and recovers below it, bounded to $[0,1]$:
\begin{equation}
\Delta f =
\begin{cases}
+\, r_\mathrm{fat}\left(\dfrac{P_\mathrm{out} - P_\mathrm{thr}}{P_\mathrm{thr}}\right)^2 \Delta t, & P_\mathrm{out} > P_\mathrm{thr}, \\[2mm]
-\, r_\mathrm{rec}\, \Delta t, & \text{otherwise.}
\end{cases}
\end{equation}
The net force advances speed and position, with speed clamped to $[0, v_\mathrm{max}]$:
\begin{equation}
F_\mathrm{net} = F_\mathrm{rider} - F_\mathrm{drag} - F_\mathrm{roll} - F_\mathrm{grav}, \qquad
v \leftarrow \mathrm{clamp}\!\Big(v + \tfrac{\Delta t}{m} F_\mathrm{net},\ 0,\ v_\mathrm{max}\Big), \qquad
x \leftarrow x + v\, \Delta t .
\end{equation}
Each segment adds its turn penalty $\tau_i$ to the clock; the finishing time is the accumulated $t$ once the full course distance is covered.

\subsection{Riders and courses}
A rider is defined by $(m, C_dA, P_\mathrm{base}, P_\mathrm{thr}, P_\mathrm{max})$; the profiles are in \texttt{data/rider\_profiles.csv}. The shared constants $C_{rr}$, $v_\mathrm{max}$, $k_\mathrm{hill}$, $r_\mathrm{fat}$, $r_\mathrm{rec}$, $c_\mathrm{fat}$ are common to all riders. A course is an ordered list of segments (distance, gradient, turn penalty, wind exposure) in \texttt{data/courses/}.

\section{Results}

Table~\ref{tab:matrix} gives the finishing time for each rider profile on each course (Tokyo is one $22.1$ km lap; the men ride two laps in the real event).
% TODO (write yourselves): interpret the matrix -- TT specialists vs climbers, the effect of terrain (flat Flanders vs rolling Tokyo), and gender.

\begin{table}[h]\centering
\caption{Finishing time (minutes) by rider profile and course, canonical model.}
\label{tab:matrix}
\begin{tabular}{lrrr}
\toprule
Rider & Tokyo (1 lap) & Flanders & Custom \\
\midrule
baseline        & 38.7 & 70.2 & 8.8 \\
male\_tt         & 31.0 & 56.5 & 7.1 \\
female\_tt       & 34.9 & 63.2 & 8.0 \\
male\_climber    & 33.5 & 62.5 & 7.8 \\
female\_climber  & 37.2 & 69.0 & 8.6 \\
\bottomrule
\end{tabular}
\end{table}

\paragraph{Still to run on the canonical model.}
% TODO (run on the notebook, then tabulate):
\begin{itemize}
  \item Wind sensitivity: vary the per-segment effective headwind (or apply a uniform offset) and report the change in finishing and split times.
  \item Power-deviation sensitivity: scale the target power by $(1 \pm \text{dev})$ and report the resulting range of split times.
  \item Validation: compare model finishing times against \texttt{data/real\_results.csv} (the men ride two Tokyo laps). The baseline rider currently runs well over the real times, so calibrate the power parameters.
\end{itemize}

\section{Model assessment}
% TODO (write yourselves): strengths and weaknesses against the problem goals; where the model is likely to fail; how it could be tested, calibrated, or improved.

\paragraph{Known limitations and possible refinements.}
\begin{itemize}
  \item Pacing is reactive: target power is a fixed function of gradient, not the time-optimal power profile under the fatigue budget.
  \item The tuning constants ($k_\mathrm{hill}$, $r_\mathrm{fat}$, $r_\mathrm{rec}$, $c_\mathrm{fat}$) are hand-calibrated, not externally sourced.
  \item Wind is treated as a scalar headwind per segment, not a direction-dependent vector.
  \item With the baseline parameters the model predicts times above the real results -- a calibration issue addressed by the literature-based rider profiles.
  \item Alternative modelling choice (required by the rubric): the critical-power (CP/$W'$) model in \texttt{archive/cp\_w\_prime\_model/}, a different fatigue framework that spends a finite anaerobic reserve $W'$ above a critical power.
\end{itemize}

\section{Conclusion}
% TODO (write yourselves): summarise the findings and give future directions.

\section*{References}
% TODO: format properly. Full list with URLs is in references/references.md and
% references/data-sources.md (course geometry, 2021 ITT results, rider power /
% CdA / Crr / air-density ranges).

\end{document}
